\section{Introduction}

Two Android packages may look alike for many legitimate and illegitimate reasons. They may be separate builds of one code base, applications that embed the same advertising library, products generated from one template, regional variants signed by different subsidiaries, siblings produced from the same prompt, or an original and an unauthorized derivative. A detector can be accurate about a shared feature while still being unable to distinguish among those histories. The central problem is therefore not merely whether two artifacts are similar. It is whether the observation is strong enough for the relationship asserted by an analyst, marketplace, developer, or incident responder.

This distinction is established indirectly across several mature literatures. Repackaging research studies code reuse, graph similarity, resource resemblance, visual clones, watermarks, and malicious grafting~\cite{droidmoss,attackclones,adrob,appink,droidmarking,piggybacking}. Third-party-library (TPL) research studies how to separate application code from reused components and identify families and versions despite shrinking, partial inclusion, and obfuscation~\cite{commonlibraries,libd,libradar,libscout,libid,atvhunter}. Software-supply-chain research instead asks who performed a step, which materials were consumed, which artifact was produced, how repository metadata was authorized, and whether different observers receive consistent histories~\cite{tuf,intoto,chainiac,sigstore,diplomat}. These bodies of work are complementary, but their guarantees do not compose automatically.

AI-mediated development makes the boundary more visible. A text-to-application system can generate compilable Android source from a natural-language description, while an agentic system can create and revise an application through multiple tool calls~\cite{text2app,appforge}. Code assistants may also introduce insecure code, propose nonexistent packages, reproduce sensitive material, or alter code through iterations whose causal history is not represented by the final APK alone~\cite{pearce,perry,sandoval,spracklen,codexleaks}. An application can therefore be novel at the byte level, similar at the behavioral or user-interface level, composed partly of inherited code, and produced through an authorized but weakly recorded process. No single binary relation captures all of those facts.

The practical consequence is a recurring category error. A high similarity score is reported as proof of copying; a component match is treated as proof of application ancestry; a valid signature is treated as proof of rightful authorization; an SBOM entry is treated as proof that the listed code is present and reachable; or a reproducible build is treated as proof that the inputs were approved. Each inference may become defensible with additional evidence, but the missing premise is often implicit. Once hidden, it is difficult to benchmark, transfer to a new workflow, or challenge during an appeal.

This survey does not claim that Android repackaging or supply-chain security lacks prior synthesis. Li, Bissyand\'e, and Klein reviewed 57 repackaging approaches and introduced RePack to expose reproducibility and comparison weaknesses~\cite{repack}. Wu et al. subsequently published a comprehensive survey of Android repackaging detection, including code- and resource-based techniques, datasets, metrics, and future directions~\cite{wurepack2026}. Zhan et al. organized 74 Android TPL studies around tasks including module separation, family identification, and version attribution, while Zeng et al. broadened the view to TPL security in application software~\cite{tpl,zengtpsurvey2024}. Recent syntheses cover software-supply-chain attacks and defenses, secure-design properties, research directions, and cyber-supply-chain risk~\cite{astra,properties,ladisa,reichert2024,williams2025,gokkaya2026,cyberrisk2026}; empirical SBOM studies additionally expose design and adoption problems in practice~\cite{bi2024sbom,bomsaway2024}. A defensible contribution must therefore be narrower than ``an updated repackaging survey'' and different from a general supply-chain taxonomy.

The question retained here is claim-centered: \emph{which observable evidence supports which Android lineage or assurance proposition, under what transformation and trust assumptions, and what bindings are required when evidence is composed across artifacts, components, build steps, release policy, behavior, and AI-mediated changes?} A bounded adversary analysis of \ReviewMatrixWorks{} review and method records found no single coded source that directly covered seven dimensions with one observation-to-claim discipline and an explicit cross-layer binding rule. This is a positioning result over the retained set, not proof that no uncoded publication contains part of the argument.

The survey answers four research questions:

\begin{description}
\item[RQ1: Claim structure.] What distinct lineage and assurance claims recur in Android analysis, and which evidence objects can support each claim?
\item[RQ2: Transformation robustness.] How do rewriting, obfuscation, resource changes, library reuse, compilation, and generation alter the validity of artifact-level evidence?
\item[RQ3: Cross-layer composition.] How can artifact analysis, component evidence, authenticated process records, release authorization, and behavior analysis be combined without overstating any layer?
\item[RQ4: AI-mediated development.] Which assumptions become fragile when source and mobile artifacts are generated or revised through nondeterministic, tool-using systems?
\end{description}

The review contributes a seven-level claim ladder; a mechanism map linking objects, transformations, oracles, ceilings, and counter-histories; a digest-bound cross-layer assurance pattern; and an analysis of AI-mediated sibling generation, synthesized dependencies, hidden context, tool state, and model drift. The ladder adds propositions rather than ranking every higher level as universally better.

The evidence base contains \CorpusWorks{} works, including \PrimaryWorks{} primary/system studies and \SecondaryWorks{} review/method sources. Every record has a claim-level extraction and inference limit; \SourceRechecks{} load-bearing assignments received same-process source rechecks. The counts are not a complete search denominator or independent dual coding.

The synthesis is that lineage evidence is \emph{non-substitutable}: artifact and component evidence recover relationships when process records are absent; authenticated records establish direction and accountability; release controls bind distribution authority; and behavior analysis addresses security properties. Composition strengthens assurance only when artifacts, principals, times, policies, and transformations are explicitly bound.

Section~\ref{sec:method} gives the review design. Sections~\ref{sec:model}--\ref{sec:ai} develop the claim model and domain syntheses; Sections~\ref{sec:synthesis} and~\ref{sec:agenda} derive the cross-layer blueprint and agenda.
